\section*{Bab \ref{ch:b3:deformation-mechanisms} --- Mekanisme Deformasi dan Zona Geser}

\begin{solution}{exo:b3:deformation-mechanisms:1}
Forsterit: $1573/2163 = 0.73$; es: $263/273 = 0.96$. Es jauh lebih dekat ke
titik lelehnya dan karena itu lebih mudah merayap: gletser mengalir di
permukaan Bumi.
\end{solution}

\begin{solution}{exo:b3:deformation-mechanisms:2}
$\gamma = \tan\ang{70} = 2.75$ dan $\tan\ang{80} = 5.67$. Di dekat \ang{90}
tangen berubah sangat cepat terhadap sudut: satu derajat galat pada $\psi =
\ang{84}$ mengubah $\gamma$ sebesar seperlima.
\end{solution}

\begin{solution}{exo:b3:deformation-mechanisms:3}
$\sigma = (10^{3.56}/10)^{1/1.26} = 363.1^{0.794} = \qty{108}{MPa}$.
\end{solution}

\begin{solution}{exo:b3:deformation-mechanisms:4}
$E/RT = 223000/(8.314\times773.15) = 34.69$, sehingga $\dot\varepsilon =
\num{1.1e-4}\times20^4\times e^{-34.69} = \num{1.1e-4}\times\num{1.6e5}\times
\num{8.6e-16} = \qty{1.5e-14}{s^{-1}}$.
\end{solution}

\begin{solution}{exo:b3:deformation-mechanisms:5}
$\eta = \sigma/(2\dot\varepsilon)$: $\num{65.4e6}/\num{2e-14} =
\qty{3.3e21}{Pa.s}$ pada \qty{400}{\celsius} dan $\num{18.0e6}/\num{2e-14} =
\qty{9.0e20}{Pa.s}$ pada \qty{500}{\celsius}: kerak tengah pada suhu ini
kira-kira sama kentalnya dengan mantel di bawah lembaran es.
\end{solution}

\begin{solution}{exo:b3:deformation-mechanisms:6}
$n = \log(\num{1.6e-4}/\num{e-5})/\log2 = \log16/\log2 = 4$: \omterm{def:b3:deformation-mechanisms:creep}{rayapan dislokasi},
dengan \omterm{def:b3:deformation-mechanisms:flow-law}{eksponen tegangan} jauh di atas 1.
\end{solution}

\begin{solution}{exo:b3:deformation-mechanisms:7}
$\ln10 = (E/R)(1/1000 - 1/1100)$, sehingga $E = 8.314\times2.3026/\num{9.091e-5} =
\qty{211}{kJ.mol^{-1}}$.
\end{solution}

\begin{solution}{exo:b3:deformation-mechanisms:8}
Nabarro--Herring ($d^{-2}$): 100 kali lebih cepat; Coble ($d^{-3}$): 1000 kali.
Sepanjang batas, atom yang berdifusi memakai jalur dengan penampang, $\delta
d$, yang menyusut seiring butir, di samping jarak yang lebih pendek dan volume
yang lebih kecil untuk dipindahkan.
\end{solution}

\begin{solution}{exo:b3:deformation-mechanisms:9}
$\sigma_b^{2.5} = \num{3.59e-6}\,d^{-3}/\num{6.54e-17}$: untuk $d = 1000$,
$\sigma_b = 54.9^{0.4} = \qty{5.0}{MPa}$; untuk $d = 100$, $\sigma_b =
\num{54900}^{0.4} = \qty{79}{MPa}$. Kemiringannya adalah $-3/2.5 = -1.2$: butir
sepersepuluh kali lebih halus menaikkan batas sebesar $10^{1.2} = 16$.
\end{solution}

\begin{solution}{exo:b3:deformation-mechanisms:10}
Pada kedua jalur, selisih konsentrasi yang menggerakkan sama dan jaraknya
adalah $d$; volume yang dipindahkan per satuan regangan adalah $d^3$. Kisi
menawarkan penampang $\propto d^2$, sehingga lajunya berbanding $d^2/(d\,d^3) =
d^{-2}$; batas menawarkan $\propto\delta d$, sehingga $\delta d/(d\,d^3) = \delta
d^{-3}$. Rasio laju batas terhadap kisi adalah $(B_B\delta D_B)/(B_LD_Ld)$,
lebih besar dari 1 ketika $d < (B_B/B_L)\,\delta D_B/D_L$.
\end{solution}

\begin{solution}{exo:b3:deformation-mechanisms:11}
$\eta \propto \dot\varepsilon^{(1-n)/n} = \dot\varepsilon^{-2.5/3.5}$: laju
sepuluh kali lebih cepat membagi viskositas dengan $10^{0.714} = 5.2$. Pita
yang kebetulan berdeformasi lebih cepat menjadi lebih lemah, sehingga ia
menanggung lebih banyak deformasi: regangan terkumpul di dalamnya.
\end{solution}

\begin{solution}{exo:b3:deformation-mechanisms:12}
Sebelum: $\dot\varepsilon_{\mathrm{dis}} = \num{6.54e-17}\times10^{3.5} =
\qty{2.07e-13}{s^{-1}}$ dan $\dot\varepsilon_{\mathrm{dif}} = \num{3.59e-6}
\times10/3000^3 = \qty{1.3e-15}{s^{-1}}$: \qty{2.08e-13}{s^{-1}}, \omterm{def:b3:deformation-mechanisms:creep}{rayapan
dislokasi}. Sesudah: $\dot\varepsilon_{\mathrm{dif}} = \num{3.59e-6}\times10/30^3 =
\qty{1.33e-9}{s^{-1}}$, \omterm{def:b3:deformation-mechanisms:creep}{rayapan difusi}, \num{6400} kali lebih cepat pada
tegangan yang sama: rekristalisasi telah melunakkan zona itu secara besar-besaran.
\end{solution}

\begin{solution}{pb:b3:deformation-mechanisms:1}
\textbf{1.} Dike berputar dari \ang{90} menjadi \ang{6} terhadap zona: $\psi =
\ang{84}$ dan $\gamma = \tan\ang{84} = 9.51$.

\textbf{2.} $\gamma w = 9.51\times400 = \qty{3804}{m}$, sekitar \qty{3.8}{km}.

\textbf{3.} $S_1 = \sqrt{1 + 9.51^2/4} + 9.51/2 = 4.859 + 4.755 = 9.614$,
$S_2 = 1/S_1 = 0.104$, dan $S_1/S_2 = 92.4$.

\textbf{4.} $\tan2\theta' = 2/9.51 = 0.210$, $\theta' = \ang{5.9}$: foliasi,
yang mendekati sumbu panjang, semakin mendekati zona seiring $\gamma$ membesar.
grows.

\textbf{5.} Jika diikuti ke utara, dike terbawa ke timur di dalam zona dan
keluar darinya lebih jauh ke timur: dinding utara bergerak ke timur relatif
terhadap dinding selatan. Dilihat dari dinding selatan, sisi jauh bergerak ke
kanan: zona itu dekstral.

\textbf{6.} Pada bidang yang memuat lineasi peregangan, horizontal di sini
untuk zona geser-mendatar, dan tegak lurus terhadap foliasi vertikal: sayatan
horizontal. Foliasi S-nya condong berlawanan arah jarum jam dari zona, ke arah
timur-timur laut, seperti sumbu panjang elips regangan.

\textbf{7.} $\sigma = (10^{3.56}/20)^{1/1.26} = 181.5^{0.794} = \qty{62}{MPa}$.

\textbf{8.} $242.1^{0.794} = \qty{78}{MPa}$ untuk \qty{15}{\micro m} dan
$145.2^{0.794} = \qty{52}{MPa}$ untuk \qty{25}{\micro m}.

\textbf{9.} Ya: eksperimen mencakup rentang \num{34} sampai \qty{268}{MPa}.

\textbf{10.} Zona itu berdeformasi pada \qty{450}{\celsius}, jauh di bawah 800 sampai
\qty{1100}{\celsius} pada eksperimen; hubungan yang tidak bergantung pada
suhu dapat dibawa turun ke sana, meskipun itu masih
merupakan ekstrapolasi.

\textbf{11.} $D = 3631/100^{1.26} = 3631/331 = \qty{11}{\micro m}$.

\textbf{12.} $E/RT = 223000/(8.314\times723.15) = 37.09$ dan $e^{-37.09} =
\num{7.79e-17}$.

\textbf{13.} $\dot\varepsilon = \num{1.1e-4}\times62^4\times\num{7.79e-17} =
\num{1.1e-4}\times\num{1.478e7}\times\num{7.79e-17} = \qty{1.27e-13}{s^{-1}}$.

\textbf{14.} $\dot\gamma = 2\dot\varepsilon = \qty{2.53e-13}{s^{-1}}$.

\textbf{15.} $\eta = \num{62e6}/(2\times\num{1.27e-13}) = \qty{2.4e20}{Pa.s}$,
mendekati \qty{3e20}{Pa.s} pada mantel atas di bawah Danau Ardo: zona itu
selemah mantel.

\textbf{16.} $e^{-E/RT}$ berubah dengan faktor 0.064 pada \qty{400}{\celsius}
(enam belas kali lebih lambat) dan 11.0 pada \qty{500}{\celsius}.

\textbf{17.} $(78/62)^4 = 2.5$ kali lebih cepat.

\textbf{18.} $t = 9.51/\num{2.53e-13} = \qty{3.75e13}{s}$.

\textbf{19.} $\num{3.75e13}/\num{3.156e10} = \qty{1190}{ka}$.

\textbf{20.} $\qty{3804}{m}/\qty{1.19e6}{yr} = \qty{3.2}{mm/yr}$.

\textbf{21.} Sepuluh kali atau lebih lambat daripada batas lempeng utama: zona itu
adalah untaian kecil, atau salah satu dari beberapa untaian yang berbagi gerakan suatu batas.

\textbf{22.} $\dot\gamma = 0.02/(\num{3.156e7}\times400) = \qty{1.58e-12}{s^{-1}}$,
6.2 kali lebih cepat; tegangan naik sebesar $6.2^{1/4} = 1.58$, menjadi \qty{98}{MPa}, dan
butirannya akan berukuran $3631/98^{1.26} = \qty{11}{\micro m}$.

\textbf{23.} Waktu berbanding terbalik dengan $\sigma^{-4}$: rentang ukuran butir saja memberi
$(62/78)^4 = 0.40$ sampai $(62/52)^4 = 2.0$ kali taksiran, dan
ketidakpastian eksponen $1.26 \pm 0.13$ dan konstanta \omterm{def:b3:deformation-mechanisms:flow-law}{hukum aliran}
memperlebarnya lebih jauh: hasilnya berupa orde besaran, bukan tanggal.

\textbf{24.} \omterm{def:b3:deformation-mechanisms:shear-zone}{Zona geser} Varden memerlukan sekitar
\textbf{\qty[text-series-to-math=true]{1190}{ka}} untuk mengakumulasi regangannya.
\end{solution}
